\documentclass[11pt]{article}
\usepackage{amsmath, amssymb}
\usepackage{geometry}
\usepackage{enumitem}
\geometry{margin=2.5cm}

\title{\textbf{Meeting Notes --- Prof.\ Asaf Cohen}\\
\large Decentralized Detection: Questions \& Insights -1}
\author{Oren Ram \& Ziv Abraham}
\date{Spring 2026}

\begin{document}
\maketitle

\begin{enumerate}[leftmargin=*, itemsep=10pt]

    \item The paper defines two cost functions: $J^N$ for finite $N$ and
    $J^N_{EE} = \frac{\log J^N}{N}$ for the infinite-sensor regime.
    What is the conceptual difference between them, and why does taking
    the logarithm give a meaningful objective?

    \item Assumption~1(ii) requires that all observation distributions
    are dominated by a common base measure $Q$.
    What is the intuition behind this condition, and what concretely
    breaks in the analysis if it does not hold?

    \item Assumption~1(iii) fixes a common action space $\mathbb{U} = \{1,\ldots,M\}$
    for all sensors, and Prof.\ Cohen's annotation suggests that
    allocating rate unevenly across sensors can improve performance.
    How exactly is that rate-allocation idea related to Assumption~1(iii),
    and is there a connection between the likelihood ratio $l^i$ at each
    sensor and the total system rate $R_{\text{total}}$?

    \item Each sensor maps its observation to a discrete action $u^i$
    and sends it to the fusion center.
    Does the model assume sensors transmit immediately upon receiving
    a signal, or do they accumulate observations and send a packet?

    \item Lemma~2 characterizes the optimal decision rule at the fusion center.
    What is the main point of the lemma, and how is it proven?

    \item With $M$ actions, the optimal sensor policy $\gamma^i$ partitions
    the likelihood-ratio axis into $M$ intervals.
    How does increasing the rate $R = \log_2 M$ affect the optimal
    policies $\gamma^i$ for sensors $i > 1$?

    \item The paper assumes sensors know the exact distributions
    $P_{Y|H_1}$ and $P_{Y|H_2}$ when designing their encoding policies.
    Can a \emph{universal} sensor policy be designed that achieves optimal
    performance without knowledge of the true distributions --- analogous
    to universal source coding?

    \item The proof relies on observations being i.i.d.\ across sensors
    given the hypothesis.
    What happens to the exchangeability and identical-policy results
    when sensors are spatially correlated?

    \item The paper considers binary hypothesis testing throughout.
    Do the main results --- threshold policies and identical encoding
    in the limit --- generalize to $M$-ary hypothesis testing?

\end{enumerate}

\end{document}